\subsection*{Chapter 3: Gravimetric Methods of Analysis}

\textbf{Relative Supersaturation (Rs):}
\begin{equation}
R_s = \frac{Q-S}{S}
\end{equation}
\textit{Q}: Concentration of the precipitating reagent at the instant of mixing, \textit{S}: Solubility of the precipitate in the medium.

\vspace{5mm}

\textbf{Calculation of Result (\% Analyte from a solid sample without dilution):}
\begin{equation}
F = \frac{M_{analyte}}{M_{sample}} (stoichiometric)
\end{equation}
\textit{Example: $F = \frac{3 \cdot FeSO_4}{Fe_3O_4}$}
\begin{equation}
\%X = m' \times F \times \frac{100}{m}
\end{equation}
\textit{\%X}: Percentage of analyte X, \textit{m'}: Mass of weighed compound (precipitate) (g), \textit{m}: Mass of sample (g), \textit{F}: Gravimetric factor.

\vspace{5mm}

\subsubsection*{Calculations for Diluted Samples}

\textbf{Solid Sample (with dilution):}
\begin{equation}
\% X = m' \times F \times \frac{V}{v} \times \frac{100}{m}
\end{equation}
\textit{V}: Total volume of the diluted sample (e.g., in mL), \textit{v}: Volume of the aliquot taken for analysis (e.g., in mL).

\vspace{5mm}

\textbf{Liquid Sample:}
\begin{equation}
X(g/L) = m' \times F \times \frac{V}{v_1} \times \frac{1000}{v}
\end{equation}
\textit{X(g/L)}: Concentration of analyte X in grams per liter, \textit{V}: Total diluted volume, \textit{v\textsubscript{1}}: Volume of aliquot taken, \textit{v}: Original volume of the liquid sample taken for dilution.

